% ── Chapter 1: Introduction ───────────────────────────────────
\chapter{Introduction}

\section{Background and Problem Statement}

Antimicrobial resistance (AMR) has emerged as one of the most pressing and complex threats to global
public health in the twenty-first century. When bacteria, viruses, fungi, or parasites evolve mechanisms
to withstand the drugs designed to eliminate them, standard treatments fail and infections persist,
spread, and increasingly result in death. The World Health Organization (WHO) has classified AMR as a
critical global priority, estimating that drug-resistant infections directly caused approximately
1.27 million deaths in 2019 alone, with indirect contributions to nearly 5 million additional
fatalities worldwide \cite{who_amr2022}. Without coordinated action, it has been projected that
AMR could claim up to 10 million lives per year by 2050, surpassing cancer as a leading cause of
mortality globally. Beyond the human toll, AMR imposes staggering economic costs through extended
hospital stays, more expensive second- and third-line therapies, and reduced workforce productivity.

At the core of combating AMR lies the practice of Antimicrobial Susceptibility Testing (AST).
Before initiating antibiotic therapy, clinicians must determine which drugs remain effective against
the specific bacterial pathogen infecting a patient. Prescribing an antibiotic to which the organism
is already resistant not only fails the patient but also selects for further resistance in the
population, perpetuating the crisis. Accurate and timely AST is therefore not merely a diagnostic
procedure---it is a public health intervention. Delays in receiving susceptibility results, which
can stretch from 24 to 72 hours using culture-based methods, force clinicians to begin empirical
broad-spectrum therapy, a practice that itself accelerates resistance development.

The Kirby-Bauer disk diffusion method, first systematically described and validated in the 1960s,
remains the most widely adopted technique for AST in clinical and reference microbiology laboratories
worldwide \cite{asm2009kirby,clsi2025}. Its enduring popularity stems from its simplicity, low cost, and
reliability when performed correctly. In this method, a Mueller-Hinton agar plate is inoculated with
a standardized suspension of the target bacterium (typically adjusted to a 0.5 McFarland turbidity
standard). Commercially manufactured paper disks, each impregnated with a precise mass of a specific
antibiotic, are placed on the inoculated agar surface at defined inter-disk distances to prevent
zone overlap. After an overnight incubation period at 35--37$^\circ$C, the antibiotic diffuses
radially outward from each disk, establishing a concentration gradient in the agar. Where the local
antibiotic concentration exceeds the minimum inhibitory concentration (MIC) of the bacterium,
bacterial growth is suppressed, resulting in a visible circular clearing called the zone of
inhibition (ZOI). Where the concentration falls below the MIC, bacterial colonies grow normally.

\begin{figure}[H]
  \centering
  \safeimg[width=0.85\textwidth]{figures/kirby_bauer_example.png}{Example of Kirby-Bauer disk diffusion plate showing zones of inhibition}
  \caption{Representative Kirby-Bauer disk diffusion plate illustrating zones of inhibition (ZOI)
    around antibiotic-impregnated disks. The diameter of each clearing zone is measured to
    determine antibiotic susceptibility.}
  \label{fig:kirby_bauer}
\end{figure}

The diameter of each ZOI is then measured and compared against published interpretive breakpoints
to classify the tested organism as Susceptible (S), Intermediate (I), or Resistant (R). These
breakpoints are rigorously maintained and updated by two major international bodies: the
Clinical and Laboratory Standards Institute (CLSI) in the United States, and the European Committee
on Antimicrobial Susceptibility Testing (EUCAST) in Europe \cite{eucast2023}. CLSI publishes
the M02 standard governing disk diffusion testing procedures, while EUCAST publishes annually
revised breakpoint tables derived from pharmacokinetic/pharmacodynamic (PK/PD) modeling,
epidemiological cutoff values (ECOFFs), and clinical outcome data \cite{giske2022}. Both
organizations provide species- and drug-specific diameter thresholds that must be applied
correctly to yield clinically meaningful results.

Despite the apparent simplicity of measuring a circle diameter, manual ZOI measurement is
considerably error-prone in practice. Laboratory technicians use rulers, Vernier calipers, or
specialized template devices to measure each zone, and inter-observer variability is a recognized
and persistent challenge in the field \cite{humphries2018}. Factors such as partially overlapping
zones between adjacent disks, heterogeneous or ``swarming'' growth within the zone boundary,
uneven agar surface, irregular plate inoculation, sub-optimal incubation, and inconsistent
photographic or visual angle all contribute to measurement variability. Even a single-millimeter
discrepancy can be clinically decisive: for example, CLSI 2025 defines a ZOI diameter of
18\,mm or above for \textit{Staphylococcus epidermidis} tested against oxacillin (1\,$\mu$g) as
Susceptible, while 17\,mm or below is classified as Resistant. With no intermediate category
for this drug-organism combination, every millimeter matters. Furthermore, laboratories processing
high case volumes require technicians to measure dozens of zones per shift, introducing fatigue
as a compounding source of error.

Existing automated AST platforms, such as the bioM\'{e}rieux VITEK 2, the BD Phoenix, and the
Microscan WalkAway systems, resolve these measurement challenges through dedicated hardware and
proprietary software. These systems offer rapid, reproducible results and are widely used in
large hospital laboratories. However, they require substantial capital investment---often tens
of thousands of US dollars per unit---and ongoing reagent and maintenance costs that are
prohibitive for smaller clinics, district hospitals, and laboratories in low- and middle-income
countries (LMICs). These resource-constrained settings are often those most severely burdened
by AMR, creating a stark disparity between the laboratories that most need accurate AST and
those that can afford automated solutions.

There is therefore a compelling case for developing low-cost, AI-driven software tools capable
of analyzing standard photographs of Kirby-Bauer AST plates. Recent research has demonstrated
that machine learning and artificial intelligence can play a significant role in predicting and
combating antimicrobial resistance \cite{bilal2025}, and advancements in AI-driven drug
sensitivity testing are rapidly expanding the practical scope of automated AST \cite{liao2025}.
Such tools could be deployed on commodity hardware---a laptop or even a smartphone---and would
require no specialized instruments beyond the agar plate itself, potentially extending the reach
of accurate AST to settings currently dependent on manual reading, including those guided by
national laboratory standards such as those maintained by the Department of Medical Sciences,
NIH Thailand \cite{dmsc_nih}. The present project addresses this need by developing, evaluating,
and deploying an AI-based system for automated ZOI detection and measurement.

\section{Motivation}

Recent advances in deep learning and computer vision have produced models that rival or exceed
human performance on a wide range of image recognition and object detection tasks. Of particular
relevance to AST image analysis is the progress in detecting and localizing circular or
elliptical structures in complex, cluttered image backgrounds. Unlike natural-image datasets where
objects vary enormously in shape, AST plates present a constrained visual domain: circular zones
on a textured agar background, with known-size reference objects (antibiotic disks) present in
every image. This domain structure is exploited by the present project's design choices.

The motivation for this project arises from two parallel observations. First, deep learning-based
object detection models---specifically two-stage detectors such as Faster R-CNN and one-stage
detectors such as YOLOv11---have demonstrated the capability to learn to identify zones of
inhibition from labeled training data, even when zones partially overlap or display atypical
morphology. Second, a comprehensive software solution for AST reading must not only detect zones
but also correctly associate each zone with its corresponding antibiotic disk. Without knowing
which antibiotic produced which zone, the diameter measurements are clinically meaningless.
This association problem requires reading the alphanumeric label printed on each antibiotic disk---a
task that falls within the domain of optical character recognition (OCR).

The combination of an object detection backbone for ZOI localization, an OCR engine for antibiotic
label reading, and a structured breakpoint database for clinical classification forms the conceptual
foundation of the system developed here. Each component presents its own technical challenges.
Object detection must cope with image noise, variable lighting, and zone overlap. OCR must read
small, low-contrast text on curved disk surfaces, potentially photographed at arbitrary orientations.
The calibration subsystem must convert pixel-scale measurements to physical millimeters using
the known physical dimensions of the antibiotic disk as an internal reference.

A further motivation is the opportunity for direct, rigorous comparison between classical
image-processing approaches and modern deep learning methods on the same dataset and evaluation
protocol. The Hough Circle Transform, a well-established technique for circular structure detection,
provides a principled baseline against which the deep learning models can be assessed. Such a
comparison provides empirical evidence regarding the practical utility of deep learning for this
specific task, beyond what is available from general-purpose benchmarks.

\section{Objectives}

This project has four primary objectives:
\begin{enumerate}
  \item To study and implement the interpretation of Kirby-Bauer disk diffusion results
        according to CLSI 2025 \cite{clsi2025} and EUCAST 2023 \cite{eucast2023} standards,
        ensuring that the system's susceptibility classifications are clinically grounded.
  \item To develop and train AI models---specifically Faster R-CNN and YOLOv11n---capable of
        detecting zones of inhibition and measuring their diameters from AST plate photographs
        with mean absolute error (MAE) at or below the clinically acceptable threshold of 3\,mm
        \cite{humphries2018}.
  \item To rigorously compare the accuracy, speed, and robustness of three detection approaches:
        Faster R-CNN (two-stage deep learning), YOLOv11n (one-stage deep learning), and Hough
        Circle Transform (classical image processing baseline), providing a quantitative basis
        for selecting the best-performing model for deployment.
  \item To deploy the system as a functional, user-facing web application suitable for practical
        use in clinical laboratory settings, integrating the detection model with the OCR pipeline
        and breakpoint classification database into a coherent, accessible tool.
\end{enumerate}

\section{Project Scope}

The scope of this project is defined and bounded as follows:

\begin{itemize}
  \item \textbf{Input data:} The system accepts photographic images of Mueller-Hinton agar plates
        used in standard Kirby-Bauer AST experiments. Images are expected to be captured from a
        top-down perspective, with consistent lighting and a fixed camera-to-plate distance to
        maintain scale consistency. Images may originate from laboratory cameras, clinical
        microscopy cameras, or standard smartphone cameras.

  \item \textbf{Detection models:} The system trains, evaluates, and compares three detection
        approaches: Faster R-CNN with a ResNet-50 feature extraction backbone, YOLOv11n
        (Ultralytics nano variant), and Hough Circle Transform implemented via OpenCV. Each
        model is evaluated on the same held-out test set under identical conditions.

  \item \textbf{Evaluation criteria:} The primary accuracy metric is mean absolute error (MAE)
        in millimeters between system-predicted ZOI diameters and ground-truth diameters measured
        by a trained microbiologist. Secondary metrics include inference speed (frames per second)
        and detection recall. The clinically accepted MAE target is below 3\,mm.

  \item \textbf{Calibration:} Pixel-to-millimeter conversion is performed using the antibiotic
        disk itself as an in-image reference, since the physical diameter of a standard antibiotic
        disk is precisely defined as 6\,mm by CLSI specifications. No external calibration target
        or ruler is required in the field of view.

  \item \textbf{Web application:} The final web application allows users to upload an image from
        file or capture one in real time via a connected camera. The system returns a dashboard
        displaying detected zones, assigned antibiotic labels (from OCR), measured diameters, and
        the final S/I/R classification for each antibiotic-disk pair based on the selected
        CLSI or EUCAST standard. Historical results are stored for audit and review.

  \item \textbf{Out of scope:} The system does not perform bacterial identification or speciation;
        the user must specify the organism being tested (e.g., \textit{Staphylococcus aureus},
        \textit{Escherichia coli}) for the correct breakpoint to be applied. The system does not
        support broth microdilution, E-test, or other AST formats. Results from this software
        tool are intended to assist, not replace, a qualified microbiologist's judgment.
\end{itemize}

\section{Expected Outcomes}

The expected tangible outputs of this project are as follows:

\begin{enumerate}
  \item \textbf{AI detection pipeline:} An end-to-end automated pipeline capable of detecting and
        measuring ZOI diameters from AST plate images. The best-performing model is expected to
        achieve an MAE below 3\,mm on the held-out test set, consistent with the clinical
        acceptability criterion established by Humphries et al.\ \cite{humphries2018}.

  \item \textbf{Comparative model analysis:} A rigorous, quantitative comparison of Faster R-CNN,
        YOLOv11n, and Hough Circle Transform on the same dataset, documenting the trade-offs
        between detection accuracy, inference speed, and implementation complexity for the AST
        domain. This analysis is intended to inform future practitioners selecting detection
        approaches for similar medical image analysis tasks.

  \item \textbf{Deployed web application:} A fully operational web application integrating the
        best-performing detection model, the EasyOCR-based antibiotic label recognition pipeline,
        and the CLSI 2025/EUCAST 2023 breakpoint classification database, providing an end-to-end
        workflow from image upload to susceptibility report.

  \item \textbf{Annotated dataset:} A curated, annotated dataset of 287 real AST plate images,
        expanded to 543 through augmentation and supplemented with 2,000 synthetic images,
        providing a reusable resource for future research in AI-assisted AST.
\end{enumerate}

\section{Report Structure}

The remainder of this report is organized as follows. Chapter~2 provides background theory and a
review of related work, covering the Kirby-Bauer AST method and its clinical standards, the
architecture and principles of the object detection models employed (Faster R-CNN, YOLOv11n, and
Hough Circle Transform), the OCR engine used for antibiotic label recognition, key image
preprocessing techniques, and a review of prior published work on automated AST image analysis.
Chapter~3 details the methodology, encompassing the full dataset collection and preparation
workflow, annotation procedures, model training protocols, calibration approach, OCR pipeline
design, susceptibility classification logic, and overall system architecture. Chapter~4 presents
experimental results including detection accuracy metrics, speed benchmarks, OCR performance, and
an analysis of error sources and calibration refinement. Chapter~5 concludes the report with a
summary of findings, limitations, and directions for future work.
